\section*{Chapitre \ref{ch:m2:the-fx-market} --- Le marché des changes}

\begin{solution}{exo:m2:the-fx-market:1}
Elle se vend au cours acheteur du dealer~: $5\,000\,000 \times 1.1462 = \text{USD}~5\,731\,000$.
La fourchette est de 2 \omterm{def:m2:the-fx-market:pair}{pips}, soit 1.74 points de base du cours médian.
\end{solution}

\begin{solution}{exo:m2:the-fx-market:2}
USDJPY, \omterm{def:m2:the-fx-market:pair}{pip} 0.01~; EURGBP, \omterm{def:m2:the-fx-market:pair}{pip} 0.0001~; AUDJPY, \omterm{def:m2:the-fx-market:pair}{pip} 0.01.
\end{solution}

\begin{solution}{exo:m2:the-fx-market:3}
Jeudi 8 octobre, sans jour férié~; vendredi 9 octobre, avec un jour férié japonais le
mercredi~; vendredi 9 octobre, avec un jour férié américain le jeudi, soit la date de valeur elle-même.
USDCAD se règle le mercredi 7 octobre.
\end{solution}

\begin{solution}{exo:m2:the-fx-market:4}
Cours acheteur $1.1462/1.3373 = 0.85710$, cours vendeur $1.1464/1.3371 = 0.85738$~: 2.78 \omterm{def:m2:the-fx-market:pair}{pips}.
\end{solution}

\begin{solution}{exo:m2:the-fx-market:5}
Chaque transaction met en jeu deux devises et chacune est comptée, de sorte que les parts de toutes
les devises s'additionnent pour donner le double du total. Le dollar figurant d'un côté de 89~\% des transactions,
un \omterm{def:m2:the-fx-market:spot}{cours croisé} comme SEKJPY est peu liquide~: l'essentiel est négocié en deux transactions en dollars, SEK
contre USD et USD contre JPY, où la liquidité est profonde.
\end{solution}

\begin{solution}{exo:m2:the-fx-market:6}
La transaction est transmise en give-up au prime broker, qui devient la contrepartie
du dealer exécutant~; le dealer fait face au risque de crédit du prime broker, non à celui du fonds. Le
prime broker fait face au fonds dans une transaction de sens inverse et en supporte le risque de crédit,
qu'il limite et encadre par des appels de marge.
\end{solution}

\begin{solution}{exo:m2:the-fx-market:7}
209.7375/209.7956, soit une fourchette relative de 2.771 points de base~; les fourchettes relatives des
jambes sont de 1.496 et 1.275, dont la somme vaut 2.771.
\end{solution}

\begin{solution}{exo:m2:the-fx-market:8}
Le volume total se répartit entre des centaines de \omterm{def:m2:the-fx-market:pair}{paires de devises}, d'instruments et de fuseaux horaires~; l'essentiel
il s'agit de swaps de change et de négociation entre dealers, et une petite devise contre une autre
peut s'échanger pour quelques millions par heure. La liquidité au prix affiché se limite à la
taille affichée~; un gros ordre est exécuté par tranches, paie une fourchette plus large ou fait bouger le prix,
et à certaines heures de la journée comme sur certaines paires, il y a très peu à négocier.
\end{solution}

\begin{solution}{pb:m2:the-fx-market:1}
\textbf{1.} Les deux cours vendeurs~: il achète des dollars contre des yens au cours vendeur USDJPY, 156.88, et
des euros contre des dollars au cours vendeur EURUSD, 1.1464.
\textbf{2.} 179.7929/179.8472.
\textbf{3.} 5.43 \omterm{def:m2:the-fx-market:pair}{pips}, 3.02 points de base.
\textbf{4.} 1.745 et 1.275 points de base, dont la somme vaut 3.02~: les fourchettes relatives s'additionnent.
\textbf{5.} USD~11\,464\,000, achetés puis immédiatement dépensés.
\textbf{6.} Oui~: son cours acheteur, 179.80, est inférieur au cours vendeur synthétique et son cours vendeur, 179.84,
supérieur au cours acheteur synthétique~; il est d'ailleurs plus serré que le synthétique.
\textbf{7.} Non~: son cours acheteur, 179.85, est supérieur au cours vendeur synthétique. Vendre des euros à B au
cours de 179.85, puis les racheter par les jambes au cours de 179.8472.
\textbf{8.} $10\,000\,000 \times (179.85 - 179.84723)$, soit JPY~27\,680, soit environ USD~176.
\textbf{9.} Un cours affiché n'a pas suivi un mouvement des jambes, à cause d'un système lent ou
d'un trader distrait~; les arbitragistes et les contrôles internes du dealer l'éliminent en
quelques instants.
\textbf{10.} Du dealer~A, à 179.84, soit en dessous du cours vendeur de B et du cours vendeur synthétique.
\textbf{11.} Deux tickets et deux fourchettes, le risque que la seconde jambe bouge
avant d'être exécutée, trois devises à régler et peut-être des tailles minimales sur les
jambes.
\textbf{12.} Lorsqu'un dealer a un flux croisé de sens inverse ou voit un intérêt dans les deux sens
pour le \omterm{def:m2:the-fx-market:spot}{cours croisé} lui-même, et qu'il peut coter à l'intérieur du synthétique.
\textbf{13.} La transaction directe échange des euros contre des yens~; la version synthétique achète et vend aussi
des dollars, qui doivent être livrés et payés, généralement via CLS.
\textbf{14.} Les jambes en yens et le \omterm{def:m2:the-fx-market:spot}{cours croisé} direct se règlent le vendredi 9 octobre~; la
jambe EURUSD se règle le jeudi 8 octobre. Les dollars achetés le jeudi sont nécessaires le
vendredi~: la voie synthétique exige un swap d'un jour pour les aligner.
\textbf{15.} Il compense le \omterm{def:m2:the-fx-market:spot}{cours croisé} avec les flux des autres clients et ne couvre que le
solde, souvent à de meilleurs prix qu'un client qui négocie les jambes.
\textbf{16.} Si les jambes bougent et que le cours affiché reste en place, il est croisé avec le
synthétique et se fait cueillir, comme celui du dealer~B.
\textbf{17.} Les paires contre dollar forment les marchés profonds~; un \omterm{def:m2:the-fx-market:spot}{cours croisé} peu liquide se cote à partir
d'elles, et son propre carnet d'ordres apporte peu.
\textbf{18.} Acheter EURJPY en gros volume par les jambes revient à acheter EURUSD et USDJPY~:
les deux montent, et le \omterm{def:m2:the-fx-market:spot}{cours croisé} synthétique aussi.
\textbf{19.} Résultat à retenir~: \emph{la borne synthétique} pour EURJPY est
179.7929/179.8472, 5.43 \omterm{def:m2:the-fx-market:pair}{pips}~: le client ne doit jamais payer plus de 179.8472 pour
des euros, puisqu'il peut les acheter par les jambes~; toute cotation directe n'est avantageuse que si
elle se situe à l'intérieur de celle-ci.
\textbf{20.} Les cours acheteur et vendeur des deux paires contre dollar, par lesquels chacun peut
négocier le \omterm{def:m2:the-fx-market:spot}{cours croisé} synthétique.
\end{solution}

\begin{solution}{iq:m2:the-fx-market:1}
Le dealer achète des euros à 1.1462 dollars et les vend à 1.1464~; la fourchette est de
deux \omterm{def:m2:the-fx-market:pair}{pips}. Un client qui achète des euros frappe le cours vendeur, en payant 1.1464 dollars pour chaque
euro.
\iqlookfor{\omterm{def:m2:the-fx-market:pair}{devise de base} et \omterm{def:m2:the-fx-market:pair}{devise de cotation}, et le côté payé par l'acheteur.}
\end{solution}

\begin{solution}{iq:m2:the-fx-market:2}
Prendre la date de transaction (après le roll de 17:00 à New York), avancer jusqu'au second jour qui
est un jour ouvré dans les deux devises, en ignorant un jour férié américain tombant le jour
intermédiaire pour les paires en dollars, et au premier pour USDCAD~; puis vérifier que la date de valeur est
ouverte dans les deux devises. Tester les week-ends, les jours fériés du jour intermédiaire et du jour de valeur
dans chaque devise, les jours fériés américains intermédiaires, les fins d'année, les transactions juste avant et
après 17:00, et les paires en T+1.
\iqlookfor{la règle, son exception pour le dollar, et une liste réelle de cas limites.}
\end{solution}

\begin{solution}{iq:m2:the-fx-market:3}
Vendre des euros contre des livres en passant par les dollars~: vendre des euros au cours acheteur EURUSD $b_1$,
acheter des livres avec ces dollars au cours vendeur GBPUSD $a_3$~: cours acheteur $b_1/a_3$. Acheter
des euros~: payer le cours vendeur EURUSD $a_1$ avec des dollars issus de la vente de livres au
cours acheteur GBPUSD $b_3$~: cours vendeur $a_1/b_3$.
\iqlookfor{les bons côtés de chaque jambe.}
\end{solution}

\begin{solution}{iq:m2:the-fx-market:4}
Ses participants sont des banques qui négocient entre elles et avec leurs clients, en
volumes importants, à travers les fuseaux horaires, dans de nombreuses devises, et règlent
bilatéralement~; aucune bourse ni aucune chambre de compensation unique n'était nécessaire ni imposée. L'électronique
a ajouté les carnets d'ordres inter-dealers, puis les plateformes multi-dealers et mono-dealer,
les algorithmes et les teneurs de marché non bancaires~; elle a réduit les coûts et fragmenté
la liquidité, tout en concentrant la négociation dans les banques qui captent le plus de flux.
\iqlookfor{l'histoire et les effets de l'électronisation.}
\end{solution}

\begin{solution}{iq:m2:the-fx-market:5}
Il permet au client de négocier avec plusieurs dealers au nom du prime broker, sur son crédit,
et il règle et déclare l'ensemble en un seul endroit~; les dealers font face au
prime broker. Il prend le risque de crédit du client~: si le client fait défaut, le prime
broker doit honorer les transactions. Il limite ce risque par des limites de négociation, des marges
et un suivi intrajournalier (\cref{ch:m2:getting-access-fx}).
\iqlookfor{le mécanisme de give-up et le risque de crédit.}
\end{solution}

\begin{solution}{iq:m2:the-fx-market:6}
Conserver le meilleur cours acheteur et le meilleur cours vendeur de chaque jambe en dollars dans un
tableau de taille fixe, mis à jour par les gestionnaires de données de marché~; à chaque mise à jour
d'une jambe, ne recalculer que les croisés qui l'utilisent, à partir d'une table de dépendances
précalculée, les règles de côté étant résolues au démarrage afin que le chemin critique
se réduise à deux multiplications ou divisions. Publier avec un numéro de séquence
et les horodatages des jambes, supprimer un croisé lorsqu'une jambe est périmée, et tester
le résultat par rapport à une implémentation de référence lente sur des données enregistrées.
\iqlookfor{dépendances précalculées, aucune allocation sur le chemin critique, contrôles de péremption.}
\end{solution}
